\block{Objetivos}{
	El propósito central de esta investigación es cuantificar la eficiencia productiva municipal mediante técnicas de aprendizaje estadístico avanzado. Los objetivos específicos incluyen:
	\begin{itemize}[leftmargin=1.5em]
		\item \textbf{Modelación Espacio-Temporal:} Desarrollar un marco probabilístico que permita capturar la heterogeneidad
		      municipal sin perder la consistencia estructural de la isla.
		\item \textbf{Descomposición de Tendencias:} Aislar los efectos de choques externos (inflación, crisis energética)
		      de la eficiencia operativa intrínseca de cada municipio mediante el uso de HSGP.
		\item \textbf{Identificación de Disparidades:} Ejecutar un "Stress Test" estadístico para identificar los
		      municipios con mayor resiliencia económica y aquellos en riesgo de estancamiento productivo.
		\item \textbf{Validación Teórica:} Contrastar las elasticidades de capital y trabajo estimadas con
		      la teoría económica clásica bajo un entorno de incertidumbre bayesiana.
	\end{itemize}
}

